\SourceLeaf{360}%
verzichtet also auf die Vortheile, die man in dem letztern in der Ausübung ungleich häufiger vorkommenden Falle bei dem Gebrauch anderer Methoden von diesem Umstande ziehen kann. Der nützliche Gebrauch der \textsc{Bessels}chen Methode wird sich also auf den Fall beschränken, wo die beiden Punkte nicht unmittelbar durch die Seite eines wirklich gemessenen Dreiecks zusammenhängen, sondern wo der Zusammenhang durch eine grössere Reihe von Dreiecken vermittelt ist. Allein dann muss man mit Recht fragen. wie denn die \textit{Data} zu der Aufgabe erlangt werden sollen, nemlich die wirkliche Länge der die beiden Punkte verbindenden geodaetischen Linie, und der Winkel, welchen sie an dem einen Endpunkte mit dem Meridian macht? Diese Bestimmung durch eine bloss \textit{sphärische} Berechnung der Übergangsdreiecke zu machen (wie \textsc{Bessel} bei der wenig ausgedehnten preussischen Gradmessung gethan hat), würde bei einer viel grössern Entfernung nicht mehr zulässig bleiben: soll aber dieser Übergang sphäroidisch gerechnet werden, so wird dies schon für sich allein eben so viel Arbeit erfordern, als wenn man gleich von jedem folgenden Punkt Breite, Länge und das rückwärts geltende Azimuth bestimmt. Übrigens gelten diese Bemerkungen auch von \textsc{Ivory}'s Auflösungsmethode, die mit der von \textsc{Bessel} viele Ähnlichkeit hat, aber das eigentliche practische Bedürfniss wenig berücksichtigt.

Über die in der vorliegenden Abhandlung gegebene Methode möge hier noch Folgendes bemerkt werden.

Die Formeln geben unmittelbar die \textit{Differenzen} zwischen den beiden Breiten und den beiden Azimuthen, so wie den Längenunterschied, und eben hierauf beruht, bei der Kleinheit dieser Differenzen (insofern rücksichtlich der Azimuthe das eine von der Südseite, das andere von der Nordseite des Meridians gezählt wird) die Schärfe der Rechnung, ohne mehrstellige Logarithmen zu erfordern. Die Symmetrie und Einfachheit der Formeln hingegen beruhet darauf, dass sie zunächst nicht von der Breite und dem Azimuthe an dem einen Endpunkte, sondern von dem Mittel der beiden Breiten und dem Mittel der beiden Azimuthe abhängen. Es folgt daraus, dass die Formeln, zur Auflösung der Aufgabe, wie sie oben ausgesprochen ist, nur vermöge eines indirecten Verfahrens oder einer successiven Annäherung benutzt werden können. Geübte und mit den Hülfen des kleinen Mechanismus derartiger Operationen vertraute Rechner werden in diesem Umstande kaum eine Unbequemlichkeit finden, zumal da man annehmen kann, dass fast immer zu der Zeit, wo die scharfe Ausführung der Rech%
